\chapter{Boolean Analysis — Threshold Functions: Basic}
%%% generated by the blueprint pipeline from TCSlib.BooleanAnalysis.ThresholdFunctions.Basic
\section{Overview}

This chapter introduces multilinear polynomial threshold representations, signed parity majorities,
and Fourier, distance, and noise quantities on the Boolean cube.
It also records elementary properties of uniform expectation and the threshold sign convention that assigns $1$ to zero.

\section{Declarations}

\begin{definition}[Walsh coefficient representation]
\lean{BooleanAnalysis.ThresholdFunctions.MultilinearPolynomial}
\leanok
For a natural number $n$, let $I=\{0,\ldots,n-1\}$ be the coordinate set.
A multilinear polynomial on the $n$-dimensional Boolean cube is represented by an arbitrary
coefficient family $p:\mathcal P(I)\to\mathbb R$, assigning a real coefficient $p(S)$ to each subset $S\subseteq I$.
\end{definition}

\begin{definition}[Evaluation in the Walsh basis]
\lean{BooleanAnalysis.ThresholdFunctions.MultilinearPolynomial.eval}
\leanok
\uses{BooleanAnalysis.BoolCube,BooleanAnalysis.ThresholdFunctions.MultilinearPolynomial,BooleanAnalysis.chiS}
Let $n\in\mathbb N$, $I=\{0,\ldots,n-1\}$, and let $p:\mathcal P(I)\to\mathbb R$ be a coefficient family.
For $x\in\{0,1\}^I$, its polynomial value is
\[
P_p(x)=\sum_{S\subseteq I}p(S)\chi_S(x),\qquad
\chi_S(x)=\prod_{i\in S}(-1)^{x_i},
\]
where the product over the empty set is $1$.
\end{definition}

\begin{definition}[Bounded multilinear degree]
\lean{BooleanAnalysis.ThresholdFunctions.MultilinearPolynomial.HasDegreeAtMost}
\leanok
\uses{BooleanAnalysis.ThresholdFunctions.MultilinearPolynomial}
Let $n,k\in\mathbb N$, $I=\{0,\ldots,n-1\}$, and $p:\mathcal P(I)\to\mathbb R$.
The coefficient family $p$ has degree at most $k$ precisely when $p(S)=0$ for every
$S\subseteq I$ with $|S|>k$.
\end{definition}

\begin{definition}[Sign with positive tie breaking]
\lean{BooleanAnalysis.ThresholdFunctions.thresholdSign}
\leanok
For every real number $r$, the threshold sign is the real number
\[
\sigma(r)=\begin{cases}1,&r\ge0,\\-1,&r<0.\end{cases}
\]
In particular, zero has threshold sign $1$.
\end{definition}

\begin{definition}[Thresholding a multilinear polynomial]
\lean{BooleanAnalysis.ThresholdFunctions.MultilinearPolynomial.threshold}
\leanok
\uses{BooleanAnalysis.BooleanFunc,BooleanAnalysis.ThresholdFunctions.MultilinearPolynomial,BooleanAnalysis.ThresholdFunctions.MultilinearPolynomial.eval,BooleanAnalysis.ThresholdFunctions.thresholdSign}
Let $n\in\mathbb N$, $I=\{0,\ldots,n-1\}$, and $p:\mathcal P(I)\to\mathbb R$.
The threshold function associated with $p$ assigns to each $x\in\{0,1\}^I$ the value
$1$ when $\sum_{S\subseteq I}p(S)\prod_{i\in S}(-1)^{x_i}\ge0$, and $-1$ otherwise.\ In particular, a zero polynomial value is assigned $1$.
\end{definition}

\begin{definition}[Polynomial sign representation]
\lean{BooleanAnalysis.ThresholdFunctions.IsPolynomialThresholdRepresentation}
\leanok
\uses{BooleanAnalysis.BooleanFunc,BooleanAnalysis.ThresholdFunctions.MultilinearPolynomial,BooleanAnalysis.ThresholdFunctions.MultilinearPolynomial.threshold}
Let $n\in\mathbb N$, $I=\{0,\ldots,n-1\}$, $f:\{0,1\}^I\to\mathbb R$, and $p:\mathcal P(I)\to\mathbb R$.
The family $p$ is a polynomial threshold representation of $f$ precisely when, for every $x\in\{0,1\}^I$,
$f(x)=\sigma(\sum_{S\subseteq I}p(S)\prod_{i\in S}(-1)^{x_i})$, where $\sigma(r)=1$ for $r\ge0$ and $\sigma(r)=-1$ for $r<0$.
\end{definition}

\begin{definition}[Degree bounded polynomial threshold functions]
\lean{BooleanAnalysis.ThresholdFunctions.IsPolynomialThreshold}
\leanok
\uses{BooleanAnalysis.BooleanFunc,BooleanAnalysis.ThresholdFunctions.IsPolynomialThresholdRepresentation,BooleanAnalysis.ThresholdFunctions.MultilinearPolynomial,BooleanAnalysis.ThresholdFunctions.MultilinearPolynomial.HasDegreeAtMost}
Let $n,k\in\mathbb N$, $I=\{0,\ldots,n-1\}$, and $f:\{0,1\}^I\to\mathbb R$.
The function $f$ is a polynomial threshold function of degree at most $k$ precisely when there exists
$p:\mathcal P(I)\to\mathbb R$ such that $p(S)=0$ whenever $|S|>k$ and
\[
f(x)=\sigma\!\left(\sum_{S\subseteq I}p(S)\prod_{i\in S}(-1)^{x_i}\right)
\quad\text{for every }x\in\{0,1\}^I,
\]
where $\sigma(r)=1$ for $r\ge0$ and $-1$ for $r<0$.
\end{definition}

\begin{definition}[Linear threshold functions]
\lean{BooleanAnalysis.ThresholdFunctions.IsLinearThreshold}
\leanok
\uses{BooleanAnalysis.BooleanFunc,BooleanAnalysis.ThresholdFunctions.IsPolynomialThreshold}
Let $n\in\mathbb N$, $I=\{0,\ldots,n-1\}$, and $f:\{0,1\}^I\to\mathbb R$.
The function $f$ is a linear threshold function precisely when it has a polynomial threshold
representation of degree at most one: there exists $p:\mathcal P(I)\to\mathbb R$ with
$p(S)=0$ for $|S|>1$ and $f(x)=\sigma(\sum_{S\subseteq I}p(S)\prod_{i\in S}(-1)^{x_i})$ for every $x$,
where $\sigma(r)=1$ for $r\ge0$ and $-1$ for $r<0$.
\end{definition}

\begin{definition}[Nonvanishing homogeneous linear representation]
\lean{BooleanAnalysis.ThresholdFunctions.IsHomogeneousLinearThresholdRepresentation}
\leanok
\uses{BooleanAnalysis.BoolCube,BooleanAnalysis.BooleanFunc,BooleanAnalysis.ThresholdFunctions.thresholdSign,BooleanAnalysis.boolToSign}
Let $n\in\mathbb N$, $I=\{0,\ldots,n-1\}$, $f:\{0,1\}^I\to\mathbb R$, and $a:I\to\mathbb R$.
The weights $a$ give a homogeneous linear threshold representation of $f$ precisely when,
for every $x\in\{0,1\}^I$, the sum $L_a(x)=\sum_{i\in I}a_i(-1)^{x_i}$ is nonzero and
$f(x)=\sigma(L_a(x))$, where $\sigma(r)=1$ for $r\ge0$ and $-1$ for $r<0$.
\end{definition}

\begin{definition}[Majorities of signed Walsh characters]
\lean{BooleanAnalysis.ThresholdFunctions.IsMajorityOfSignedParities}
\leanok
\uses{BooleanAnalysis.BooleanFunc,BooleanAnalysis.ThresholdFunctions.thresholdSign,BooleanAnalysis.chiS}
Let $n,s\in\mathbb N$, $I=\{0,\ldots,n-1\}$, and $f:\{0,1\}^I\to\mathbb R$.
The function $f$ is a majority of $s$ signed parities precisely when there are subsets
$S_0,\ldots,S_{s-1}\subseteq I$ and signs $\varepsilon_j\in\{-1,1\}$ such that
\[
f(x)=\sigma\!\left(\sum_{j=0}^{s-1}\varepsilon_j\prod_{i\in S_j}(-1)^{x_i}\right)
\quad\text{for every }x\in\{0,1\}^I,
\]
where repetitions of the subsets are allowed and $\sigma(r)=1$ for $r\ge0$, $-1$ for $r<0$;
thus ties, including the empty sum when $s=0$, give the value $1$.
\end{definition}

\begin{definition}[Nonzero coefficient support]
\lean{BooleanAnalysis.ThresholdFunctions.MultilinearPolynomial.support}
\leanok
\uses{BooleanAnalysis.ThresholdFunctions.MultilinearPolynomial}
Let $n\in\mathbb N$, $I=\{0,\ldots,n-1\}$, and $p:\mathcal P(I)\to\mathbb R$.
The support of $p$ is the finite set $\operatorname{supp}(p)=\{S\subseteq I:p(S)\ne0\}$.
\end{definition}

\begin{definition}[Number of nonzero coefficients]
\lean{BooleanAnalysis.ThresholdFunctions.MultilinearPolynomial.sparsity}
\leanok
\uses{BooleanAnalysis.ThresholdFunctions.MultilinearPolynomial,BooleanAnalysis.ThresholdFunctions.MultilinearPolynomial.support}
Let $n\in\mathbb N$, $I=\{0,\ldots,n-1\}$, and $p:\mathcal P(I)\to\mathbb R$.
The sparsity of $p$ is $|\{S\subseteq I:p(S)\ne0\}|$, the number of its nonzero coefficients.
\end{definition}

\begin{definition}[Fourier mass through a level]
\lean{BooleanAnalysis.ThresholdFunctions.fourierWeightUpTo}
\leanok
\uses{BooleanAnalysis.BooleanFunc,BooleanAnalysis.fourierCoeff}
Let $n,k\in\mathbb N$, $I=\{0,\ldots,n-1\}$, and $f:\{0,1\}^I\to\mathbb R$.
Its Fourier weight on levels at most $k$ is
\[
W^{\le k}(f)=\sum_{\substack{S\subseteq I\\|S|\le k}}\widehat f(S)^2,
\qquad \widehat f(S)=2^{-n}\sum_{x\in\{0,1\}^I}f(x)\prod_{i\in S}(-1)^{x_i}.
\]
\end{definition}

\begin{definition}[Fourier mass beyond a level]
\lean{BooleanAnalysis.ThresholdFunctions.fourierWeightAbove}
\leanok
\uses{BooleanAnalysis.BooleanFunc,BooleanAnalysis.fourierCoeff}
Let $n,k\in\mathbb N$, $I=\{0,\ldots,n-1\}$, and $f:\{0,1\}^I\to\mathbb R$.
Its Fourier weight on levels strictly above $k$ is
\[
W^{>k}(f)=\sum_{\substack{S\subseteq I\\ k<|S|}}\widehat f(S)^2,
\qquad \widehat f(S)=2^{-n}\sum_{x\in\{0,1\}^I}f(x)\prod_{i\in S}(-1)^{x_i}.
\]
\end{definition}

\begin{definition}[Sum of absolute Fourier coefficients]
\lean{BooleanAnalysis.ThresholdFunctions.spectralOneNorm}
\leanok
\uses{BooleanAnalysis.BooleanFunc,BooleanAnalysis.fourierCoeff}
Let $n\in\mathbb N$, $I=\{0,\ldots,n-1\}$, and $f:\{0,1\}^I\to\mathbb R$.
Its spectral one-norm is
\[
\|\widehat f\|_1=\sum_{S\subseteq I}|\widehat f(S)|,
\qquad \widehat f(S)=2^{-n}\sum_{x\in\{0,1\}^I}f(x)\prod_{i\in S}(-1)^{x_i}.
\]
\end{definition}

\begin{definition}[Linear part of the Fourier expansion]
\lean{BooleanAnalysis.ThresholdFunctions.degreeOnePart}
\leanok
\uses{BooleanAnalysis.BooleanFunc,BooleanAnalysis.chiS,BooleanAnalysis.fourierCoeff}
Let $n\in\mathbb N$, $I=\{0,\ldots,n-1\}$, and $f:\{0,1\}^I\to\mathbb R$.
The degree-one part of $f$ is the function
\[
f^{=1}(x)=\sum_{i\in I}\widehat f(\{i\})(-1)^{x_i},
\qquad \widehat f(\{i\})=2^{-n}\sum_{y\in\{0,1\}^I}f(y)(-1)^{y_i}.
\]
The constant Fourier coefficient is excluded.
\end{definition}

\begin{definition}[Largest absolute coordinate influence]
\lean{BooleanAnalysis.ThresholdFunctions.maxInfluence}
\leanok
\uses{BooleanAnalysis.BooleanFunc,BooleanAnalysis.influence}
Let $n\in\mathbb N$, $I=\{0,\ldots,n-1\}$, and $f:\{0,1\}^I\to\mathbb R$.
The maximum influence is
\[
M(f)=\sup_{i\in I}|\operatorname{Inf}_i(f)|,
\qquad \operatorname{Inf}_i(f)=2^{-n}\sum_{x\in\{0,1\}^I}
\frac{(f(x)-f(x^{\oplus i}))^2}{4},
\]
where $x^{\oplus i}$ is obtained by flipping coordinate $i$, and the supremum over the empty
coordinate set is assigned $0$ when $n=0$.
\end{definition}

\begin{definition}[Uniform probability of disagreement]
\lean{BooleanAnalysis.ThresholdFunctions.disagreementProbability}
\leanok
\uses{BooleanAnalysis.BooleanFunc,BooleanAnalysis.expect}
Let $n\in\mathbb N$ and let $f,g:\{0,1\}^n\to\mathbb R$ be arbitrary real-valued functions.
Their disagreement probability is
$d(f,g)=2^{-n}|\{x\in\{0,1\}^n:f(x)\ne g(x)\}|$, using the uniform distribution on the cube.
\end{definition}

\begin{definition}[Closeness in disagreement probability]
\lean{BooleanAnalysis.ThresholdFunctions.IsClose}
\leanok
\uses{BooleanAnalysis.BooleanFunc,BooleanAnalysis.ThresholdFunctions.disagreementProbability}
Let $n\in\mathbb N$, $f,g:\{0,1\}^n\to\mathbb R$, and $\varepsilon\in\mathbb R$.
The functions $f$ and $g$ are $\varepsilon$-close precisely when
$2^{-n}|\{x\in\{0,1\}^n:f(x)\ne g(x)\}|\le\varepsilon$.
\end{definition}

\begin{definition}[Maximum absolute function value]
\lean{BooleanAnalysis.ThresholdFunctions.supNorm}
\leanok
\uses{BooleanAnalysis.BoolCube,BooleanAnalysis.BooleanFunc}
Let $n\in\mathbb N$ and $f:\{0,1\}^n\to\mathbb R$.
The supremum norm of $f$ is $\|f\|_\infty=\sup_{x\in\{0,1\}^n}|f(x)|$, equivalently
the maximum of $|f|$ on this finite nonempty cube.
\end{definition}

\begin{definition}[Stability under Fourier noise]
\lean{BooleanAnalysis.ThresholdFunctions.noiseStability}
\leanok
\uses{BooleanAnalysis.BooleanFunc,BooleanAnalysis.innerProduct,BooleanAnalysis.noiseOp}
Let $n\in\mathbb N$, $I=\{0,\ldots,n-1\}$, $f:\{0,1\}^I\to\mathbb R$, and $\rho\in\mathbb R$.
The noise stability at $\rho$ is
\[
\operatorname{Stab}_\rho(f)=2^{-n}\sum_{x\in\{0,1\}^I} f(x)(T_\rho f)(x),
\qquad (T_\rho f)(x)=\sum_{S\subseteq I}\rho^{|S|}\widehat f(S)\chi_S(x),
\]
where $\chi_S(x)=\prod_{i\in S}(-1)^{x_i}$ and
$\widehat f(S)=2^{-n}\sum_{x\in\{0,1\}^I}f(x)\chi_S(x)$; this quantity is defined for every real $\rho$.
\end{definition}

\begin{definition}[Sensitivity at a real noise parameter]
\lean{BooleanAnalysis.ThresholdFunctions.noiseSensitivity}
\leanok
\uses{BooleanAnalysis.BooleanFunc,BooleanAnalysis.ThresholdFunctions.noiseStability}
Let $n\in\mathbb N$, $I=\{0,\ldots,n-1\}$, $f:\{0,1\}^I\to\mathbb R$, and $\delta\in\mathbb R$.
The noise sensitivity is the real number
\[
\operatorname{NS}_\delta(f)=\frac{1-2^{-n}\sum_{x\in\{0,1\}^I}
f(x)\sum_{S\subseteq I}(1-2\delta)^{|S|}\widehat f(S)\chi_S(x)}{2},
\]
where $\chi_S(x)=\prod_{i\in S}(-1)^{x_i}$ and
$\widehat f(S)=2^{-n}\sum_{x\in\{0,1\}^I}f(x)\chi_S(x)$; neither $\delta\in[0,1]$
nor $f(x)\in\{-1,1\}$ is required.
\end{definition}

\begin{definition}[Parity of a binary inner product]
\lean{BooleanAnalysis.ThresholdFunctions.innerProductModTwo}
\leanok
\uses{BooleanAnalysis.BooleanFunc}
For $n\in\mathbb N$, a point of $\{0,1\}^{n+n}$ is written as two blocks
$(u,v)$ with $u,v\in\{0,1\}^n$.
The inner-product-mod-two function is
$\operatorname{IP}_n(u,v)=\prod_{i=0}^{n-1}(-1)^{u_i v_i}=(-1)^{\sum_{i=0}^{n-1}u_i v_i}$,
so it is $-1$ exactly when the number of coordinates where both bits are $1$ is odd;
the empty product for $n=0$ is $1$.
\end{definition}

\begin{definition}[Indicator of prescribed coordinates]
\lean{BooleanAnalysis.ThresholdFunctions.subcubeIndicator}
\leanok
\uses{BooleanAnalysis.BooleanFunc}
Let $n\in\mathbb N$, $I=\{0,\ldots,n-1\}$, $J\subseteq I$, and $b\in\{0,1\}^I$.
The subcube indicator associated with $J$ and $b$ is the function taking value $1$ at
$x\in\{0,1\}^I$ if $x_i=b_i$ for every $i\in J$, and value $0$ otherwise.
\end{definition}

\begin{definition}[Indicator of a normalized Hamming halfspace]
\lean{BooleanAnalysis.ThresholdFunctions.hammingBallIndicator}
\leanok
\uses{BooleanAnalysis.BooleanFunc,BooleanAnalysis.boolToSign}
Let $n\in\mathbb N$ and $t\in\mathbb R$.
The normalized Hamming-ball indicator takes value $1$ at $x\in\{0,1\}^n$ precisely when
$t\le (\sum_{i=0}^{n-1}(-1)^{x_i})/\sqrt n$, and value $0$ otherwise, including the boundary in the value-$1$ region.
When $n=0$, the quotient is assigned $0$, so the function takes value $1$ exactly when $t\le0$.
\end{definition}

\begin{lemma}[Monotonicity of uniform averaging]
\lean{BooleanAnalysis.ThresholdFunctions.expect_mono}
\leanok
\uses{BooleanAnalysis.BooleanFunc,BooleanAnalysis.expect,BooleanAnalysis.expect_eq_fintypeExpect}
\difficulty{3}
Let $n\in\mathbb N$ and $f,g:\{0,1\}^n\to\mathbb R$ satisfy $f(x)\le g(x)$ for every $x$.
Then $\mathbb E[f]\le\mathbb E[g]$, where $\mathbb E[h]=2^{-n}\sum_{x\in\{0,1\}^n}h(x)$ is uniform expectation.
\end{lemma}

\begin{lemma}[Nonnegative uniform averages]
\lean{BooleanAnalysis.ThresholdFunctions.expect_nonneg}
\leanok
\uses{BooleanAnalysis.BooleanFunc,BooleanAnalysis.expect,BooleanAnalysis.expect_eq_fintypeExpect}
\difficulty{3}
Let $n\in\mathbb N$ and $f:\{0,1\}^n\to\mathbb R$ satisfy $f(x)\ge0$ for every $x$.
Then $\mathbb E[f]\ge0$, where $\mathbb E[f]=2^{-n}\sum_{x\in\{0,1\}^n}f(x)$.
\end{lemma}

\begin{lemma}[Average of a constant]
\lean{BooleanAnalysis.ThresholdFunctions.expect_const}
\leanok
\uses{BooleanAnalysis.BoolCube,BooleanAnalysis.expect,BooleanAnalysis.expect_eq_fintypeExpect}
\difficulty{3}
For every $n\in\mathbb N$ and $c\in\mathbb R$, the uniform expectation of the constant
function $x\mapsto c$ on $\{0,1\}^n$ equals $c$:
$2^{-n}\sum_{x\in\{0,1\}^n}c=c$.
\end{lemma}

\begin{lemma}[Additivity of uniform averaging]
\lean{BooleanAnalysis.ThresholdFunctions.expect_add}
\leanok
\uses{BooleanAnalysis.BooleanFunc,BooleanAnalysis.expect}
\difficulty{1}
Let $n\in\mathbb N$ and $f,g:\{0,1\}^n\to\mathbb R$.
Then $\mathbb E[f+g]=\mathbb E[f]+\mathbb E[g]$, where addition of functions is pointwise
and $\mathbb E[h]=2^{-n}\sum_{x\in\{0,1\}^n}h(x)$.
\end{lemma}

\begin{lemma}[Scalars commute with uniform averaging]
\lean{BooleanAnalysis.ThresholdFunctions.expect_const_mul}
\leanok
\uses{BooleanAnalysis.BooleanFunc,BooleanAnalysis.expect}
\difficulty{3}
Let $n\in\mathbb N$, $c\in\mathbb R$, and $f:\{0,1\}^n\to\mathbb R$.
Then $\mathbb E[cf]=c\mathbb E[f]$, where $(cf)(x)=cf(x)$ and
$\mathbb E[h]=2^{-n}\sum_{x\in\{0,1\}^n}h(x)$.
\end{lemma}

\begin{lemma}[Negation commutes with uniform averaging]
\lean{BooleanAnalysis.ThresholdFunctions.expect_neg}
\leanok
\uses{BooleanAnalysis.BooleanFunc,BooleanAnalysis.ThresholdFunctions.expect_const_mul,BooleanAnalysis.expect}
\difficulty{1}
Let $n\in\mathbb N$ and $f:\{0,1\}^n\to\mathbb R$.
Then $\mathbb E[-f]=-\mathbb E[f]$, where $(-f)(x)=-f(x)$ and
$\mathbb E[h]=2^{-n}\sum_{x\in\{0,1\}^n}h(x)$.
\end{lemma}

\begin{lemma}[Subtraction commutes with uniform averaging]
\lean{BooleanAnalysis.ThresholdFunctions.expect_sub}
\leanok
\uses{BooleanAnalysis.BooleanFunc,BooleanAnalysis.ThresholdFunctions.expect_add,BooleanAnalysis.ThresholdFunctions.expect_neg,BooleanAnalysis.expect}
\difficulty{1}
Let $n\in\mathbb N$ and $f,g:\{0,1\}^n\to\mathbb R$.
Then $\mathbb E[f-g]=\mathbb E[f]-\mathbb E[g]$, where $(f-g)(x)=f(x)-g(x)$ and
$\mathbb E[h]=2^{-n}\sum_{x\in\{0,1\}^n}h(x)$.
\end{lemma}

\begin{lemma}[Zero average forces pointwise vanishing]
\lean{BooleanAnalysis.ThresholdFunctions.eq_zero_of_expect_eq_zero}
\leanok
\uses{BooleanAnalysis.BooleanFunc,BooleanAnalysis.expect,BooleanAnalysis.expect_eq_fintypeExpect}
\difficulty{3}
Let $n\in\mathbb N$ and $f:\{0,1\}^n\to\mathbb R$ satisfy $f(x)\ge0$ for every $x$.
If its uniform expectation $2^{-n}\sum_{x\in\{0,1\}^n}f(x)$ is zero, then $f(x)=0$
for every $x\in\{0,1\}^n$.
\end{lemma}

\begin{lemma}[Sign on the nonnegative half line]
\lean{BooleanAnalysis.ThresholdFunctions.thresholdSign_of_nonneg}
\leanok
\uses{BooleanAnalysis.ThresholdFunctions.thresholdSign}
\difficulty{1}
Let $r\in\mathbb R$ satisfy $r\ge0$, and define $\sigma(t)=1$ for $t\ge0$ and
$\sigma(t)=-1$ for $t<0$.
Then $\sigma(r)=1$, including when $r=0$.
\end{lemma}

\begin{lemma}[Sign on the negative half line]
\lean{BooleanAnalysis.ThresholdFunctions.thresholdSign_of_neg}
\leanok
\uses{BooleanAnalysis.ThresholdFunctions.thresholdSign}
\difficulty{1}
Let $r\in\mathbb R$ satisfy $r<0$, and define $\sigma(t)=1$ for $t\ge0$ and
$\sigma(t)=-1$ for $t<0$.
Then $\sigma(r)=-1$.
\end{lemma}

\begin{lemma}[Threshold signs are unit signs]
\lean{BooleanAnalysis.ThresholdFunctions.thresholdSign_pm_one}
\leanok
\uses{BooleanAnalysis.ThresholdFunctions.thresholdSign}
\difficulty{4}
For every $r\in\mathbb R$, let $\sigma(r)=1$ when $r\ge0$ and $\sigma(r)=-1$ when $r<0$.
Then $\sigma(r)=1$ or $\sigma(r)=-1$.
\end{lemma}

\begin{lemma}[Absolute size of a threshold sign]
\lean{BooleanAnalysis.ThresholdFunctions.abs_thresholdSign}
\leanok
\uses{BooleanAnalysis.ThresholdFunctions.thresholdSign,BooleanAnalysis.ThresholdFunctions.thresholdSign_pm_one}
\difficulty{4}
For every $r\in\mathbb R$, define $\sigma(r)=1$ when $r\ge0$ and $\sigma(r)=-1$ when $r<0$.
Then $|\sigma(r)|=1$.
\end{lemma}

\begin{lemma}[Multiplication by sign gives absolute value]
\lean{BooleanAnalysis.ThresholdFunctions.mul_thresholdSign_self}
\leanok
\uses{BooleanAnalysis.ThresholdFunctions.thresholdSign}
\difficulty{4}
For every $r\in\mathbb R$, define $\sigma(r)=1$ when $r\ge0$ and $\sigma(r)=-1$ when $r<0$.
Then $r\sigma(r)=|r|$.
\end{lemma}

\begin{lemma}[Recovering a number from magnitude and sign]
\lean{BooleanAnalysis.ThresholdFunctions.abs_mul_thresholdSign}
\leanok
\uses{BooleanAnalysis.ThresholdFunctions.thresholdSign}
\difficulty{4}
For every $r\in\mathbb R$, define $\sigma(r)=1$ when $r\ge0$ and $\sigma(r)=-1$ when $r<0$.
Then $|r|\sigma(r)=r$.
\end{lemma}

\begin{lemma}[Order preservation by threshold signs]
\lean{BooleanAnalysis.ThresholdFunctions.thresholdSign_mono}
\leanok
\uses{BooleanAnalysis.ThresholdFunctions.thresholdSign,BooleanAnalysis.ThresholdFunctions.thresholdSign_of_neg,BooleanAnalysis.ThresholdFunctions.thresholdSign_of_nonneg,BooleanAnalysis.ThresholdFunctions.thresholdSign_pm_one}
\difficulty{5}
Let $r,s\in\mathbb R$ satisfy $r\le s$, and define $\sigma(t)=1$ when $t\ge0$ and
$\sigma(t)=-1$ when $t<0$.
Then $\sigma(r)\le\sigma(s)$.
\end{lemma}

\begin{lemma}[Positive scaling preserves threshold signs]
\lean{BooleanAnalysis.ThresholdFunctions.thresholdSign_mul_of_pos}
\leanok
\uses{BooleanAnalysis.ThresholdFunctions.thresholdSign}
\difficulty{4}
Let $c,r\in\mathbb R$ with $c>0$, and define $\sigma(t)=1$ when $t\ge0$ and
$\sigma(t)=-1$ when $t<0$.
Then $\sigma(cr)=\sigma(r)$.
\end{lemma}

\begin{lemma}[A nearby number has the prescribed sign]
\lean{BooleanAnalysis.ThresholdFunctions.thresholdSign_eq_of_abs_sub_lt_one}
\leanok
\uses{BooleanAnalysis.ThresholdFunctions.thresholdSign,BooleanAnalysis.ThresholdFunctions.thresholdSign_of_neg,BooleanAnalysis.ThresholdFunctions.thresholdSign_of_nonneg}
\difficulty{4}
Let $y,r\in\mathbb R$ with $y\in\{-1,1\}$ and $|y-r|<1$.
For $\sigma(t)=1$ when $t\ge0$ and $\sigma(t)=-1$ when $t<0$, one has $\sigma(r)=y$.
\end{lemma}

\begin{lemma}[Pointwise thresholds take unit sign values]
\lean{BooleanAnalysis.ThresholdFunctions.isPmOne_thresholdSign}
\leanok
\uses{BooleanAnalysis.BoolCube,BooleanAnalysis.ThresholdFunctions.thresholdSign,BooleanAnalysis.ThresholdFunctions.thresholdSign_pm_one,BooleanAnalysis.isPmOne}
\difficulty{1}
Let $n\in\mathbb N$ and $g:\{0,1\}^n\to\mathbb R$, and define
$\sigma(t)=1$ when $t\ge0$ and $\sigma(t)=-1$ when $t<0$.
Then the function $x\mapsto\sigma(g(x))$ takes values in $\{-1,1\}$ at every point of the cube.
\end{lemma}

\begin{lemma}[Nonnegativity of spectral absolute mass]
\lean{BooleanAnalysis.ThresholdFunctions.spectralOneNorm_nonneg}
\leanok
\uses{BooleanAnalysis.BooleanFunc,BooleanAnalysis.ThresholdFunctions.spectralOneNorm}
\difficulty{1}
Let $n\in\mathbb N$, $I=\{0,\ldots,n-1\}$, and $f:\{0,1\}^I\to\mathbb R$.
Then its spectral one-norm is nonnegative:
\[
\sum_{S\subseteq I}|\widehat f(S)|\ge0,
\qquad \widehat f(S)=2^{-n}\sum_{x\in\{0,1\}^I}f(x)\prod_{i\in S}(-1)^{x_i}.
\]
\end{lemma}
